% SHARD-ID: AQM-46-AFFINE-LINE-SYMMETRY-WEYL-NET
% SHARD-TITLE: Affine-Line Symmetries and Weyl-Net Automorphisms
% SHARD-SUMMARY: Proves that the base-preserving automorphism group of \(\mathbf A^1_{\mathbb F_q}\) is \(x\mapsto ax+b\).
% SHARD-SUMMARY: Defines the induced action on closed points, residue fields, finite Weyl labels, and local algebras.
% SHARD-SUMMARY: Proves the resulting covariant Weyl-net automorphism and finite-support Clifford implementation.
% SHARD-KEYWORDS: affine line, automorphism group, finite field, Weyl net, Clifford, symplectic map

\section{Affine-Line Symmetries and Weyl-Net Automorphisms}
\label{sec:affine-line-symmetry-weyl-net}

\subsection{Status and Source Anchors}

This shard extends AQM-29.  It fixes the word ``symmetry'' to mean a
base-preserving \(k\)-scheme automorphism of
\(\mathbf A^1_k=\Spec k[x]\), where \(k=\mathbb F_q\).  The affine-space
definition is source-local from the Stacks Project snapshot
\path{references/algebraic_geometry/stacks_project_constructions.tex}, lines
733--758.  The affine-scheme anti-equivalence and morphism-to-affine theorem
are in \path{references/algebraic_geometry/stacks_project_schemes.tex},
lines 841--963.  Spectrum, \(D(f)\), and functoriality are in
\path{references/algebraic_geometry/stacks_project_algebra.tex}, lines
2972--2983 and 3123--3160.  The finite Weyl operators use
Ashikhmin--Knill,
\path{references/symplectic_qecc/ashikhmin_knill_0005008_source/n9.tex},
lines 249--313.  Odd-characteristic Clifford language is grounded by Gross,
\path{references/symplectic_qecc/gross_0602001_source/poswig.tex}, lines
340--537.

\subsection{Lamport Dependency Skeleton}

\[
\begin{array}{rcl}
D1&:&k=\mathbb F_q,\quad X=\mathbf A^1_k=\Spec k[x].\\
D2&:&\operatorname{Aut}_k(X)\text{ means automorphisms over }\Spec k.\\
L1&:&\operatorname{Aut}_k(X)=
      \{g_{a,b}:x\mapsto ax+b,\ a\in k^\times,\ b\in k\}.\\
D3&:&x_\pi=(\pi),\quad K_\pi=k[x]/(\pi),\quad
      \pi\text{ monic irreducible}.\\
L2&:&g_{a,b}(x_\pi)=x_{\pi^g},\quad
      \pi^g(X)=a^{\deg\pi}\pi((X-b)/a).\\
D4&:&\tau_{g,\pi}:K_{\pi^g}\to K_\pi,\quad
      \tau_{g,\pi}([X])=a[X]+b.\\
L3&:&\tau_{g,\pi}\text{ is a }k\text{-field isomorphism and preserves
      absolute traces.}\\
D5&:&E(U)=\bigoplus_{x_\pi\in U}K_\pi^2,\quad
      \mathcal A(U)=\operatorname{span}\{W_U(e):e\in E(U)\}.\\
D6&:&S_g:E(U)\to E(gU)\text{ is residue pushforward by }\tau^{-1}.\\
T1&:&S_g\text{ preserves the Weyl multiplier and commutator}.\\
T2&:&\alpha_g^U(W_U(e))=W_{gU}(S_g e)\text{ is a covariant net
      }*\text{-isomorphism}.\\
T3&:&\text{On finite supports, }\alpha_g\text{ has a tensor-permutation
      Clifford implementation.}
\end{array}
\]

\subsection{The Base-Preserving Symmetry Group D1--L1}

Let \(k=\mathbb F_q\).  By the Stacks affine-space source,
\[
  \mathbf A^1_k=\Spec k[x].
\]
A base-preserving symmetry means an automorphism of this affine scheme over
\(\Spec k\).  By the affine-scheme anti-equivalence, this is the same data as
a \(k\)-algebra automorphism
\[
  \theta:k[x]\longrightarrow k[x].
\]

\paragraph{Lemma \(L1\).}
Every \(k\)-scheme automorphism of \(\mathbf A^1_k\) is uniquely
\[
  g_{a,b}:x\longmapsto ax+b,
  \qquad a\in k^\times,\ b\in k.
\]
The group law is
\[
  g_{a,b}g_{c,d}=g_{ac,ad+b}.
\]

\emph{Proof.}
Let \(\theta:k[x]\to k[x]\) be a \(k\)-algebra automorphism.  Put
\(h=\theta(x)\).  Since \(k[x]\) is generated as a \(k\)-algebra by \(x\),
\(\theta\) is determined by \(h\).  Let \(\psi=\theta^{-1}\) and put
\(r=\psi(x)\).  Then
\[
  x=(\theta\psi)(x)=\theta(r)=r(h).
\]
If \(\deg h=m\) and \(\deg r=n\), then \(\deg r(h)=mn\).  Since
\(\deg x=1\), we have \(mn=1\).  Thus \(m=n=1\), so \(h=ax+b\) with
\(a\ne0\).  Conversely, every polynomial \(ax+b\) with \(a\ne0\) gives the
inverse \(x\mapsto a^{-1}(x-b)\).  Composition of functions gives the
displayed law.

This proves the user's proposed group exactly in the base-preserving sense.
Absolute automorphisms that are allowed to move the coefficient field are a
different question and are not used in this shard.

\subsection{Closed Points and Residue Fields D3--L3}

By AQM-29, the closed points of \(X\) are
\[
  x_\pi=(\pi),\qquad \pi\in k[x]\text{ monic irreducible}.
\]
The residue field at \(x_\pi\) is
\[
  K_\pi=\kappa(x_\pi)=k[x]/(\pi).
\]
Let \(g=g_{a,b}\).  Its coordinate pullback is
\[
  \theta_g(f)(x)=f(ax+b).
\]

\paragraph{Lemma \(L2\).}
If \(\pi\) is monic irreducible of degree \(d\), then
\[
  g(x_\pi)=x_{\pi^g},\qquad
  \pi^g(X)=a^d\pi((X-b)/a).
\]
The polynomial \(\pi^g\) is monic and irreducible.

\emph{Proof.}
Let \(\alpha=[x]\in K_\pi\).  The point image is the prime
\(\theta_g^{-1}((\pi))\).  A polynomial \(f\) lies in that prime exactly when
\[
  f(a\alpha+b)=0\quad\text{in }K_\pi.
\]
Thus the image closed point is the minimal polynomial over \(k\) of
\(\beta=a\alpha+b\).  The displayed \(\pi^g\) is monic, because the leading
coefficient is \(a^d a^{-d}=1\), and \(\pi^g(\beta)=0\).  Since
\(\alpha=a^{-1}(\beta-b)\), the fields \(k(\alpha)\) and \(k(\beta)\) are
equal.  Hence \(\beta\) has degree \(d\), so \(\pi^g\) is its irreducible
minimal polynomial.

\paragraph{Definition \(D4\).}
The residue pullback at \(x_\pi\) is
\[
  \tau_{g,\pi}:K_{\pi^g}\longrightarrow K_\pi,\qquad
  \tau_{g,\pi}([X])=a[X]+b.
\]

\paragraph{Lemma \(L3\).}
\(\tau_{g,\pi}\) is a \(k\)-field isomorphism.  For every
\(u\in K_{\pi^g}\),
\[
  \operatorname{Tr}_{K_{\pi^g}/\mathbb F_p}(u)
  =
  \operatorname{Tr}_{K_\pi/\mathbb F_p}(\tau_{g,\pi}u).
\]

\emph{Proof.}
The relation \(\pi^g(X)=0\) is sent to
\(\pi^g(a\alpha+b)=0\), so the displayed formula defines a homomorphism.
It is onto because \(\alpha=a^{-1}(\tau([X])-b)\) lies in its image.  The two
finite fields have the same size \(q^d\), hence the map is an isomorphism.
As a \(k\)-linear isomorphism, it conjugates the \(k\)-linear multiplication
map \(m_u\) on \(K_{\pi^g}\) to \(m_{\tau u}\) on \(K_\pi\).  Conjugate
linear maps have the same trace over \(k\); composing with
\(\operatorname{Tr}_{k/\mathbb F_p}\) gives the absolute-trace identity.

\subsection{The Weyl Label Pushforward D5--T1}

For an open \(U\subset X\), define
\[
  E(U)=\bigoplus_{x_\pi\in U}K_\pi^2
\]
with finite support.  A label is written
\[
  e=(q_\pi,p_\pi)_\pi .
\]
At \(x_\pi\), use
\[
  T_\pi(q)|s\rangle=|s+q\rangle,\qquad
  R_\pi(p)|s\rangle=
  \exp\!\left(\frac{2\pi i}{p_0}
  \operatorname{Tr}_{K_\pi/\mathbb F_{p_0}}(ps)\right)|s\rangle,
\]
where \(p_0=\operatorname{char}k\), and \(W_\pi(q,p)=T_\pi(q)R_\pi(p)\).
For finite support, \(W_U(e)\) is the tensor product of the local operators.

\paragraph{Definition \(D6\).}
The pushforward
\[
  S_g:E(U)\longrightarrow E(gU)
\]
is defined by
\[
  (S_g e)_{\pi^g}
  =
  \bigl(\tau_{g,\pi}^{-1}(q_\pi),
        \tau_{g,\pi}^{-1}(p_\pi)\bigr).
\]

\paragraph{Theorem \(T1\).}
\(S_g\) preserves the Weyl multiplier and commutator.  In particular,
\[
  W_{gU}(S_ge)W_{gU}(S_gf)
\]
has the same scalar multiplier as \(W_U(e)W_U(f)\).

\emph{Proof.}
It is enough to check one closed point.  Let
\((Q,P)=\tau^{-1}(q,p)\) and \((Q',P')=\tau^{-1}(q',p')\).  By Lemma \(L3\),
\[
  \operatorname{Tr}_{K_{\pi^g}/\mathbb F_{p_0}}(PQ')
  =
  \operatorname{Tr}_{K_\pi/\mathbb F_{p_0}}(pq').
\]
The same identity with \(P,Q\) and \(P',Q'\) exchanged gives preservation of
\[
  \operatorname{Tr}(PQ'-QP').
\]
These are exactly the local multiplier and commutator exponents.  Multiplying
over the finite support proves the claim.

\subsection{Net Automorphism and Clifford Implementation T2--T3}

\paragraph{Theorem \(T2\).}
For every \(g\in\operatorname{Aut}_k(X)\) and open \(U\), the rule
\[
  \alpha_g^U(W_U(e))=W_{gU}(S_ge)
\]
extends uniquely to a unital \(*\)-isomorphism
\[
  \alpha_g^U:\mathcal A(U)\longrightarrow\mathcal A(gU).
\]
These maps are covariant for inclusions and satisfy
\[
  \alpha_g^{hU}\alpha_h^U=\alpha_{gh}^U.
\]

\emph{Proof.}
The Weyl operators form a basis on every finite support, and
\(\mathcal A(U)\) is the union of the finite-support algebras.  Theorem \(T1\)
shows that the assignment respects products and adjoints on each finite
support, hence on the union.  It is bijective because \(S_{g^{-1}}\) is the
inverse label map.  If \(U\subset V\), extension by zero commutes with
pushforward of finite supports, so the square
\[
  \mathcal A(U)\to\mathcal A(V),\qquad
  \mathcal A(gU)\to\mathcal A(gV)
\]
commutes.  The residue maps \(\tau\) compose because the coordinate pullbacks
\(\theta_g\) compose; therefore \(S_gS_h=S_{gh}\), and the displayed group law
follows on the Weyl basis.

\paragraph{Theorem \(T3\).}
For a finite closed support \(S\), the finite-volume map
\[
  \alpha_g^S:\mathcal A(S)\to\mathcal A(gS)
\]
is implemented by a unitary tensor relabelling.  If
\[
  \mathcal H(S)=\bigotimes_{x_\pi\in S}\ell^2(K_\pi),
\]
define
\[
  U_g^S\left(\bigotimes_{x_\pi\in S}|s_\pi\rangle\right)
  =
  \bigotimes_{x_{\pi^g}\in gS}
  |\tau_{g,\pi}^{-1}(s_\pi)\rangle .
\]
Then
\[
  U_g^S W_S(e)(U_g^S)^*=W_{gS}(S_ge).
\]

\emph{Proof.}
At one site, residue relabelling sends
\(|s\rangle\) to \(|\tau^{-1}s\rangle\).  It carries translation by \(q\) to
translation by \(\tau^{-1}q\).  It carries the phase
\(\exp(2\pi i\,\operatorname{Tr}(ps)/p_0)\) to the phase with
\(\tau^{-1}p\) and \(\tau^{-1}s\), by Lemma \(L3\).  Tensoring the one-site
calculation over \(S\) proves the formula.

Thus the affine-line symmetry acts on the Weyl local net by symplectic
finite-support label maps.  In odd characteristic, Gross's Clifford structure
theorem makes each finite-volume implementation a Clifford unitary.  In even
characteristic, the algebra automorphism above is still proved, but the
half-Weyl Clifford convention is not imported.
